\begin{abstract}
Lattice rules are commonly constructed to perform well at sample sizes that
are powers of a fixed base. In applications, however, the available budget
may stop between these preferred sizes. We consider component-by-component
(CBC) constructions that optimize an aggregate of the expected squared
discrepancies of all prefixes up to a chosen horizon. The aggregate uses
sample-size weights distinct from the coordinate weights of the function
space. For base-two lattice sequences, dyadic difference multiplicities
reduce the criterion to a weighted kernel sum. A reversed digit-permutation
scrambling, called Newo, leaves every prefix's shift-averaged squared
discrepancy unchanged when applied to the scalar van der Corput sequence
before multiplication by a fixed lattice generator. We present a direct
permutation proof and its implication for CBC construction. For Kronecker
sequences, a lag-based expression yields a related CBC objective. Historical
comparisons from the special-session talk motivate the construction; the
numerical study and the newer Kronecker algorithm remain under development.
\keywords{quasi-Monte Carlo, lattice sequence, Kronecker sequence,
component-by-component construction, arbitrary sample size, scrambling.}
\end{abstract}

\par\smallskip
\noindent\codex{[Revision later: Revise the abstract after the newer Kronecker
construction and final experiments are in place. The present abstract describes
a discussion draft, not the final submission's completed contributions.]}
\par\smallskip
